\documentclass[10pt]{article}

\usepackage[utf8]{inputenc}

\usepackage[T1]{fontenc}

\usepackage{amsmath}

\usepackage{amsfonts}

\usepackage{amssymb}

\usepackage[version=4]{mhchem}

\usepackage{stmaryrd}

\usepackage{bbold}



\begin{document}

от $N$-мерных линейных многообразий, т. е. характеризует аппроксимативные возможности $N$-мерных линейных многообразий. Другие равносилыны и общепринятые определения $d_{N}$ таковы (см. [1]):



\[

\begin{gathered}

d_{N}(C, X)=\inf _{\left\{M_{N}\right\}} \sup _{x \in C} \inf _{y \in M_{N}}\|x-y\|= \\

=\inf \inf \left\{\varepsilon>0 \mid C \subset M_{N}+\varepsilon B\right\} .

\end{gathered}

\]



Если $\mathfrak{U}=\mathfrak{M}_{N}$ (или $=\mathfrak{R}_{N}-$ совокупности всех подпространств $\left\{L_{N}\right\}$ размерности $\leqslant N$, а $F\left(C, M_{N}\right)(F(C$, $\left.L_{N}\right)$ ) - совокупность всех аффинных (линейных) непрерывных отображений из $C$ в $M_{N}\left(L_{N}\right)$, то величина $(1)$, обозначаемая $\alpha_{N}(C, X)\left(\lambda_{N}(C, X)\right)$ и называемая а ффинды м (лине й ным) $N$-п о пе р е чви ко м, характеризует аппроксимативные возможности афффинных (линейных) $N$-мерных отображений.



Если $\mathfrak{U}=\mathfrak{\Re}_{N}$ есть совокупность $\left\{K_{N}\right\}$ всех $N$-мерных компактов (или, равнозначно, всех $N$-мерных полиэдров), а $F\left(C, K_{N}\right)$ - множество всех непрерывных отображений из $C$ в $K_{N}$, то величина (1), называемая $N$-п опе р е чни ко м п о А ле к с а нд р о в у и обозначаемая $a_{N}(C, X)$, характеризует степень аппроксимации множества $C N$-мерными компактами.



Если $\mathfrak{U}=\Sigma_{N}$ - совокупность $\left\{\xi_{N}\right\}$ всех $N$-точечных множеств $\xi_{N}\left\{x_{1}, \ldots, x_{N}\right\}$ в $X$, а $F\left(C, \xi_{N}\right)$ - совокупность всех отображений из $C$ в $\xi_{N}$, то П. (1), обозначаемый $\varepsilon_{N}(C, X)$, характеризует наименьпее уклонение данного $C$ от $N$-точечных множеств, т. е. характеризует аппроксимативные возможности $N$-точечных множеств.



Все введенные выше апироксимативные П. зависят от объемлющего пространства $X$ и могут изменяться при погружении $C$ с өго метрикой в другое нормированное пространство.



Другой тип П. связан с задачами «кодирования» элементов множества $C$ элементами другой природы или, как этот процесс еще иначе называют, с задачей восстановления. Пусть $C$ - метрич. пространство, $Z=$ $=\{\zeta\}-$ нек-рая совок упность «кодирующих" множеств, $\varphi(C, \zeta)$ - нек-рая совокупность отображений $f: C \rightarrow$ $\rightarrow \zeta$. Наконец, $\Phi=\Phi(C, Z)=\bigcup_{\zeta \in Z} \varphi(C, \zeta)-$ заданная совокупность методов кодирования. Величина



\[

p^{\Phi}(C)=\inf _{f \in \Phi} \sup _{z \in f(C)} D\left(f^{-1}(z)\right),

\]



где через $D(E)$ обозначен диаметр множества $E$, характеризует восстановимость элементов множества $C$ по информации, «закодированной» әлементами множеств $\zeta$ из $Z$ с помощью отображений из Ф. Большинство П., связанных с процессами восстановления, задается по типу (2).



Если Ф - совокупность всевозможных отображений из $C$ в $Z$, состоящего из одного множества $\{1, \ldots, N\}$, то П. (2), обозначаемый $\varepsilon^{N}(C)$, характеризует точность восстановления элемента с помощью таблицы, состоящей из $N$ элементов. Если $C$ лежит в линейном нормированном пространстве $X$ и $Ф-$ совокупность всех непрерывных аффинных отображений из $X$ в $\mathbb{R} N$, то величина, равная $p^{\Phi}(C) / 2$, где $p^{\Phi}(C)$ определено в (2), называемая $N$ - п о п е р еч н и о м по ф а н д у и обозначаемая $d^{N}(C)$, характеризует точность восстановления әлементов по их образам при аффинных отображениях в $\mathbb{R N}$. Для центрально-симметричных множеств величины $d^{N}$ имеется другое равносильное определение:



$d^{N}(C)=\inf \sup \|x\|=\inf \sup \{\varepsilon>0 \mid C \cap L \subset \varepsilon B\},\left(2^{\prime}\right)$ $\{N\} \quad x \in C \cap L N \quad\{N\}^{\varepsilon}$



где $L^{N}$ - замкнутое подпространство коразмерности $N$. Пусть $Z$ состоит из всех $N$-мерных компактов $\left\{K_{N}\right\}$, $\varphi\left(C, K_{N}\right)$ - совокупность всех непрерывных отобра- жений из $C$ в $K_{N}, \Phi=\bigcup_{\left\{K_{N}\right\}} \varphi_{(}\left(C, K_{N}\right)$, тогда (2) называется $N$ - по пе р е ч ик о мпо У р ы с он у и обозначается $u_{N}(C)$. Другое равносильное и общепринятое определение поперечника Урысона таково: $u_{N}(C)$ есть нижняя грань диаметров покрытий множества $C$ кратности $\leqslant N+1$. Поперечник по Урысону характеризует степень $N$-мерности (с точки зрения брауәровской размерности) множества $C$.



Впервые величину, названную впоследствии П., ввел в 1923 П. С. Урысон (см. [2]), когда он определил $u_{N}$. В 1933 П. С. Александров [3] вскрыл аппроксимативные аспекты теории размерности, что привело его к определению $a_{N}$. В 1936 А. Н. Колмогоров [1] определил $d_{N}-$ именно этот П. наиболее интенсивно изучался далее в теории приближений. В 1931 Л. С. Понтрягин и Л. Г. Шнирельман (см. [4]) выразили размерность (топологич. характеристику), использовав асимптотическую метрич. характеристику $N_{\varepsilon}(C)$ (обратную к поперечнику $\left.\varepsilon^{N}(C)\right)$, равную для метрич. пространства $C$ наименьшему числу элементов $2 \varepsilon-$ покрытия для множества $C$. Интерес к подобным величинам возрос в 50-х гг. когда А. Н. Колмогоров [5], базируясь на идеях теории информации, ввел величину $N_{\varepsilon}(C, X)$ (обратную к поперечнику $\left.\varepsilon_{N}(C, X)\right)$ и сформулировал программу исследований величин типа $N_{\varepsilon}(C), N_{\varepsilon}(C, X)$ и им подобных как специальный раздел теории приближений, связанный с вопросами наилучшего табулирования функций. Двоичный логарифмм величины $N_{\varepsilon}(C, X)$ получил название $\varepsilon$-энтропии множества $C$, а $\log _{2} N_{\varepsilon}(C)-$ а 6 с ол ю т н о й $\varepsilon-$ э н т р о п и и множества $C$.



Получено множество конкретных результатов, где те или иные П. (названные выше и иные) вычислялись для различных функциональных классов и геометрич. объектов. Такие вычисления можно разделить на две группы - асимптотические и точные.



Вот нек-рые результаты, касающиеся асимптотич. вычислений П. соболевских классов. Пусть $W_{p}^{r}-$ совокупность функций $x(\cdot)$ на конечном отрезке (скажем, на $[0,1])$, у к-рых $(r-1)$-я произвотная абсолютно непрерывна и для $r$-й производной выполнено неравенство



$\left\|x^{(r)}(\cdot)\right\|_{p}==\left(\int_{0}^{1}\left|x^{(r)}(t)\right|^{p} d t\right)^{1 / p} \leqslant 1, p \geqslant 1, k=1,2, \ldots$,



Доказана следуюпцая асимптотич. формула:



$d_{N}\left(W_{p}^{r}, L_{q}\right) \asymp\left\{\begin{array}{l}N^{-\left(r-(1 / p-1 / q)_{+}\right)}, \\ 1 \leqslant q \leqslant p \leqslant \infty \text { или } 1 \leqslant p \leqslant q \leqslant 2, \\ \left.N^{-(r-(1 / p-1 / 2)}{ }_{+}\right), \\ 1 \leqslant p \leqslant q \leqslant \infty, q \geqslant 2 .\end{array}\right.$



Из частных случаев верхней строки формулы (3) следует, что асимптотически наплучшим аппроксимирующим пространством является пространство тригонометрич. полиномов или пространство силайнов с равномерно распределенпыми узлами.



Ожидалось, что всегда имеет место такая асимптотика, т. е. подпространство тригонометрич. полиномов данной степени всегда будет асимптотически экстремальным. Однако оказалось, что это не так (см. [10], [13]). Тригонометрич. полиномы $\operatorname{lin}\{\cos n t, \sin n t, 0 \leqslant$ $\leqslant n \leqslant N\}$ оказались асимптотически не әкстремальными. Однако в ряде случаев «переставленные» гармоники, т. е. $N$ гармоник, взятых в «неправильном» порядке, все-таки оказались экстремальными.



Решение задачи о П. соболевских классов опирается на исследование вопроса о поперечнике $n$-октаэдров B $\mathbb{R} n$ :



\[

O_{p}^{n}=\left\{\left.x \in \mathbb{R}^{n}\left|\sum_{i=1}^{n}\right| x_{i}\right|^{p} \leqslant 1\right\} .

\]





\end{document}